\[
\textbf{Proof.}
\]

Fix \(x\in(0,1)\). Throughout, \(F_x\) and \(F_\psi\) denote the partial derivatives of \(F\) with respect to its first and second arguments.

For Part \(B\), it is convenient to suppress the fixed \(x\)-variable and write
\[
F_x(\psi):=F_x(x,\psi),\qquad F_\psi(\psi):=F_\psi(x,\psi),\qquad w(\psi):=w(x,\psi).
\]

From the assumptions
\[
F(x,-\psi)=F(x,\psi),\qquad F(x,\psi+\pi)=F(x,\psi),
\]
we may differentiate to obtain
\[
F_x(-\psi)=F_x(\psi),\qquad F_x(\psi+\pi)=F_x(\psi),
\]
and
\[
F_\psi(-\psi)=-F_\psi(\psi),\qquad F_\psi(\psi+\pi)=F_\psi(\psi).
\]
In particular,
\[
F_x\!\left(\psi+\frac{\pi}{2}\right)=F_x\!\left(\psi-\frac{\pi}{2}\right),
\qquad
F_\psi\!\left(\psi+\frac{\pi}{2}\right)=F_\psi\!\left(\psi-\frac{\pi}{2}\right),
\]
because adding or subtracting \(\pi\) does not change either derivative. Also,
\[
w\!\left(\psi+\frac{\pi}{2}\right)=w(\psi),
\]
so \(w\) is also \(2\pi\)-periodic.

---

### Part \(A\)

By \(F(x,\psi)=G_A(q(x,\psi))\) and the coordinate formula from \(SP2(iv)\),
\[
\partial_\phi
=
-\sqrt{1-x^2}\sin\psi\,\partial_x
+
\frac{x\cos\psi}{\sqrt{1-x^2}}\,\partial_\psi,
\]
we get
\[
\begin{aligned}
(\partial_\phi G_A)(q(x,\psi))
&=
\left(
-\sqrt{1-x^2}\sin\psi\,\partial_x
+
\frac{x\cos\psi}{\sqrt{1-x^2}}\,\partial_\psi
\right)F(x,\psi)\\
&=
-\sqrt{1-x^2}\sin\psi\,F_x(x,\psi)
+
\frac{x\cos\psi}{\sqrt{1-x^2}}\,F_\psi(x,\psi).
\end{aligned}
\]

For \(G_B\), we use
\[
G_B(q(x,\psi))=F\!\left(x,\psi-\frac{\pi}{2}\right).
\]
Hence
\[
\begin{aligned}
(\partial_\phi G_B)(q(x,\psi))
&=
\left(
-\sqrt{1-x^2}\sin\psi\,\partial_x
+
\frac{x\cos\psi}{\sqrt{1-x^2}}\,\partial_\psi
\right)
F\!\left(x,\psi-\frac{\pi}{2}\right)\\
&=
-\sqrt{1-x^2}\sin\psi\,F_x\!\left(x,\psi-\frac{\pi}{2}\right)
+
\frac{x\cos\psi}{\sqrt{1-x^2}}\,F_\psi\!\left(x,\psi-\frac{\pi}{2}\right),
\end{aligned}
\]
since
\[
\partial_x F\!\left(x,\psi-\frac{\pi}{2}\right)=F_x\!\left(x,\psi-\frac{\pi}{2}\right),
\qquad
\partial_\psi F\!\left(x,\psi-\frac{\pi}{2}\right)=F_\psi\!\left(x,\psi-\frac{\pi}{2}\right).
\]

This is exactly the desired expansion.

---

### Part \(B\)

Let
\[
I_1:=\int_0^{2\pi}\sin\psi\cos\psi\,F_x(\psi)\,F_\psi\!\left(\psi-\frac{\pi}{2}\right)\,w(\psi)\,d\psi,
\]
and
\[
I_2:=\int_0^{2\pi}\sin\psi\cos\psi\,F_\psi(\psi)\,F_x\!\left(\psi-\frac{\pi}{2}\right)\,w(\psi)\,d\psi.
\]
Then the cross-term integral to be shown equal to \(0\) is exactly \(I_1+I_2\).

We compute \(I_1\) by the change of variables
\[
\psi=u+\frac{\pi}{2}.
\]
Then \(d\psi=du\), and
\[
\psi-\frac{\pi}{2}=u.
\]
So
\[
\begin{aligned}
I_1
&=
\int_{\pi/2}^{5\pi/2}
\sin\!\left(u+\frac{\pi}{2}\right)\cos\!\left(u+\frac{\pi}{2}\right)
F_x\!\left(u+\frac{\pi}{2}\right)
F_\psi(u)\,
w\!\left(u+\frac{\pi}{2}\right)\,du.
\end{aligned}
\]
Using
\[
\sin\!\left(u+\frac{\pi}{2}\right)=\cos u,\qquad
\cos\!\left(u+\frac{\pi}{2}\right)=-\sin u,
\qquad
w\!\left(u+\frac{\pi}{2}\right)=w(u),
\]
this becomes
\[
I_1
=
-\int_{\pi/2}^{5\pi/2}
\sin u\cos u\,
F_x\!\left(u+\frac{\pi}{2}\right)
F_\psi(u)\,
w(u)\,du.
\]
The integrand is \(2\pi\)-periodic in \(u\), so we may replace the interval \([\pi/2,5\pi/2]\) by \([0,2\pi]\):
\[
I_1
=
-\int_0^{2\pi}
\sin u\cos u\,
F_x\!\left(u+\frac{\pi}{2}\right)
F_\psi(u)\,
w(u)\,du.
\]
Now use the \(\pi\)-periodicity of \(F_x\):
\[
F_x\!\left(u+\frac{\pi}{2}\right)=F_x\!\left(u-\frac{\pi}{2}\right).
\]
Therefore
\[
\begin{aligned}
I_1
&=
-\int_0^{2\pi}
\sin u\cos u\,
F_\psi(u)\,
F_x\!\left(u-\frac{\pi}{2}\right)
w(u)\,du\\
&=-I_2.
\end{aligned}
\]
Hence
\[
I_1+I_2=0.
\]

This is exactly
\[
\int_0^{2\pi}\sin\psi\cos\psi
\left[
F_x(x,\psi)\,F_\psi\!\left(x,\psi-\frac{\pi}{2}\right)
+
F_\psi(x,\psi)\,F_x\!\left(x,\psi-\frac{\pi}{2}\right)
\right]
w(x,\psi)\,d\psi
=0.
\]

---

### Part \(C\)

Now multiply the two formulas from Part \(A\). Using the shorthand from Part \(B\),
\[
\begin{aligned}
(\partial_\phi G_A)(q(x,\psi))(\partial_\phi G_B)(q(x,\psi))
&=
\Bigl(
-\sqrt{1-x^2}\sin\psi\,F_x(\psi)
+
\frac{x\cos\psi}{\sqrt{1-x^2}}\,F_\psi(\psi)
\Bigr)\\
&\quad\times
\Bigl(
-\sqrt{1-x^2}\sin\psi\,F_x\!\left(\psi-\frac{\pi}{2}\right)
+
\frac{x\cos\psi}{\sqrt{1-x^2}}\,F_\psi\!\left(\psi-\frac{\pi}{2}\right)
\Bigr).
\end{aligned}
\]
Expanding gives
\[
\begin{aligned}
(\partial_\phi G_A)(q(x,\psi))(\partial_\phi G_B)(q(x,\psi))
&=
(1-x^2)\sin^2\psi\,F_x(\psi)\,F_x\!\left(\psi-\frac{\pi}{2}\right)\\
&\quad
-x\sin\psi\cos\psi
\left[
F_x(\psi)\,F_\psi\!\left(\psi-\frac{\pi}{2}\right)
+
F_\psi(\psi)\,F_x\!\left(\psi-\frac{\pi}{2}\right)
\right]\\
&\quad
+\frac{x^2\cos^2\psi}{1-x^2}\,
F_\psi(\psi)\,F_\psi\!\left(\psi-\frac{\pi}{2}\right).
\end{aligned}
\]
Integrating against \(w(\psi)\,d\psi\) over \([0,2\pi]\), the middle term vanishes by Part \(B\). Therefore
\[
\begin{aligned}
J(x)
&=
\int_0^{2\pi}
\Biggl[
(1-x^2)\sin^2\psi\,F_x(\psi)\,F_x\!\left(\psi-\frac{\pi}{2}\right)
+
\frac{x^2\cos^2\psi}{1-x^2}\,
F_\psi(\psi)\,F_\psi\!\left(\psi-\frac{\pi}{2}\right)
\Biggr]
w(\psi)\,d\psi.
\end{aligned}
\]
Restoring the suppressed \(x\)-dependence, this is
\[
J(x)
=
\int_0^{2\pi}
\left[
(1-x^2)\sin^2\psi\,F_x(x,\psi)\,F_x\!\left(x,\psi-\frac{\pi}{2}\right)
+
\frac{x^2\cos^2\psi}{1-x^2}\,
F_\psi(x,\psi)\,F_\psi\!\left(x,\psi-\frac{\pi}{2}\right)
\right]
w(x,\psi)\,d\psi.
\]

This proves Part \(C\), and hence all three parts.